\section{The quantitative affine bound}\label{q:affine-transfer}

The progression estimate for Liouville gives the same logarithmic
exponent for every fixed affine pair. The finite initial interval
affects only the constant.

\begin{proof}[Proof of \Cref{thm:q-affine}]
Fix $a_1,a_2\ge1$ and $b_1,b_2\ge0$ with nonzero determinant, and
put
\[
 l=a_1a_2,\qquad c_0=\min(a_2b_1,a_1b_2),\qquad
 h=|a_1b_2-a_2b_1|>0.
\]
Multiplying the two affine arguments by $a_2$ and $a_1$, respectively,
and using complete multiplicativity gives
\begin{equation}\label{q:quant-affine-identity}
 \lambda(a_1n+b_1)\lambda(a_2n+b_2)
 =\lambda(l)\lambda(ln+c_0)\lambda(ln+c_0+h).
\end{equation}
Here $\lambda(l)^2=1$. For real $Y\ge0$ set
\[
 P(Y)=\sum_{1\le m\le Y}
       \lambda(m)\lambda(m+h)\1_{m\equiv c_0\pmod l}.
\]
The sum is empty for $Y<1$. Let $c>0$ be the absolute exponent in
\Cref{q:progression}. That estimate applies to the residue $c_0$
without any coprimality restriction and gives
\[
 |P(Y)|\ll_{l,h,c_0}\frac{Y}{(\log Y)^c}\qquad(Y\ge3).
\]
For every real $X\ge3$, the integers in the class $c_0\pmod l$
with $c_0<m\le lX+c_0$ are precisely $m=ln+c_0$ with
$1\le n\le X$. Consequently \eqref{q:quant-affine-identity}
gives the exact endpoint identity
\begin{equation}\label{q:affine-rate-endpoints}
 \sum_{1\le n\le X}\lambda(a_1n+b_1)\lambda(a_2n+b_2)
   =\lambda(l)\bigl(P(lX+c_0)-P(c_0)\bigr).
\end{equation}
This identity includes both determinant signs and noninteger cutoffs.
Since $lX+c_0\ge X\ge3$ and
$lX+c_0\le(l+c_0/3)X$, the first term satisfies
\[
 |P(lX+c_0)|
 \ll_{a_1,a_2,b_1,b_2}\frac{X}{(\log X)^c}.
\]
Also $|P(c_0)|\le c_0$, which is absorbed into the same bound:
the function $(\log X)^c/X$ is bounded on $[3,\infty)$.
The exponent has not changed and is independent of the affine
coefficients; only the implied constant depends on them. This proves
\Cref{thm:q-affine} for all real $X\ge3$.
\end{proof}
